\documentclass[10pt,a4paper]{amsart}
\usepackage[margin=19mm]{geometry}
\usepackage{amsmath,amssymb,mathtools}
\usepackage[scr=rsfs,cal=euler]{mathalfa}
\usepackage{xfrac}
\usepackage[unicode,hidelinks,bookmarks=false]{hyperref}
\newcommand{\ie}{i.e.\ }
\newcommand{\cf}{cf.\ }
\newcommand{\resp}{respectively\ }
\newcommand{\Subsection}{\subsection}
\providecommand{\nobreakdash}{\nobreak\hbox{-}}
\setlength{\emergencystretch}{2em}
\multlinegap0pt
\usepackage{bm}
\usepackage{bbm}
\usepackage{dsfont}
\usepackage{xspace}
\usepackage{cancel}
\usepackage{mathtools}
\usepackage{amsthm}
\makeatletter
\@ifundefined{defn}{\newtheorem{defn}{Defn}}{}
\@ifundefined{proposition}{\newtheorem{proposition}[defn]{Proposition}}{}
\@ifundefined{remark}{\newtheorem{remark}[defn]{Remark}}{}
\makeatother
\makeatletter
\newcommand{\OO}{\mathcal{O}}
\newcommand{\Q}{\mathbb Q}
\expandafter\ifx\csname alphenum\endcsname\relax
\newenvironment{alphenum}{\hfill \begin{enumerate} \renewcommand{\theenumi}{\alph{enumi}}}{\end{enumerate}}
\fi
\expandafter\ifx\csname romanenum\endcsname\relax
\newenvironment{romanenum}{\hfill \begin{enumerate} \renewcommand{\theenumi}{\roman{enumi}}}{\end{enumerate}}
\fi
\makeatother
\title[Models of CM elliptic curves with a prescribed $\ell$-adic G]{Models of CM elliptic curves with a prescribed $\ell$-adic Galois image}
\author{Enrique Gonz\'alez-Jim\'enez, \'Alvaro Lozano-Robledo,, Benjamin York}
\date{Source: arXiv:2408.04159v3 (CC BY 4.0)}
\begin{document}
\maketitle

\noindent\emph{Source and licence.} Models of CM elliptic curves with a prescribed $\ell$-adic Galois image. Version arXiv:2408.04159v3 (CC BY 4.0), distributed under the Creative Commons Attribution 4.0 International licence, as recorded in the arXiv licence metadata for this exact version and shown on the abstract page https://arxiv.org/abs/2408.04159v3. Full rights notice: licenses/arxiv-2408.04159-CC-BY-4.0.txt.\par\smallskip

\noindent\textit{Editorial scope.} Counted unit: the Proposition whose source label is \texttt{prop-8divfieldm8} in the article (source0.tex lines 918--930, source label \texttt{prop-8divfieldm8}), delivered together with the article's own complete original proof of it (source0.tex lines 932--968, one \texttt{proof} environment of 37 lines). No mathematics is written, repaired or re-derived: statement and proof are the article's own text line for line. Required context retained uncounted: rem-howtocomputedivisionfield (lines 913--916); prop-8divfieldm16 (lines 1019--1030). Every definition, display and label of those lines is retained unchanged. Omitted from this package and not redistributed: 1--912, 917--917, 931--931, 969--1018, 1031--2324. Nothing inside a retained excerpt is omitted or altered: every retained body is delivered exactly as the article's lines are written, so no omission receipt of any kind accompanies this package. Citations: the retained excerpts carry no citation command at all, so no bibliography entry of the article is transcribed and no reference is redistributed. Reference adaptation: none was needed and none was made; every reference inside the delivered unit resolves inside it, so no unresolved reference or citation is produced. Printed slips kept literal: no printed word, symbol, hypothesis or number of the counted statement or of its proof was altered. Layout and class: the article's own class is \texttt{amsart} with packages including amsfonts, epsfig, latexsym, amssymb, amsmath, amsthm, graphics, xy, fontenc, multirow, hhline, array, booktabs, ctable; the wrapper is the lane's pinned \texttt{amsart} editorial wrapper at 10pt on A4, which restates the theorem-like environments the retained text needs and adds the article's own macro definitions that the retained text needs (\texttt{\textbackslash OO}, \texttt{\textbackslash Q}), with extra wrapper packages (bm, bbm, dsfont, xspace, cancel, mathtools). The delivered numbering is the unit's own. Assets: no retained span contains a figure environment, a graphics-inclusion command, a PDF-page inclusion, a table or a TikZ picture; no image file is copied into this package, the package declares no asset dependency and no image file is redistributed with it. Licence: version arXiv:2408.04159v3 of this article is distributed under the Creative Commons Attribution 4.0 International licence, as recorded in the arXiv licence metadata for this exact version and shown on the abstract page https://arxiv.org/abs/2408.04159v3. A full rights notice is delivered at licenses/arxiv-2408.04159-CC-BY-4.0.txt. Characters: every retained excerpt is pure ASCII, so the delivered engine, XeLaTeX with its default Latin Modern text font, reports no missing character.
\begin{remark}\label{rem-howtocomputedivisionfield}
    Let $E: y^2=f(x)$ be an elliptic curve over $\Q$, let $d\in\Q^\ast$ be a square-free rational number, and let $E^d : dy^2=f(x)$ be the quadratic twist of $E$ by $d$. There is an isomorphism $\phi\colon E\to E^d$ given by $(x,y)\mapsto (x,y/\sqrt{d})$ defined over $\Q(\sqrt{d})$, and therefore, for any $N\geq 1$, we have an equality of division fields $\Q(x(E[N]))=\Q(x(E^d[N]))$. Moreover, suppose that $E[N]=\langle (x_1,y_1), (x_2,y_2) \rangle$. Then, $E^d[N]=\langle (x_1,y_1/\sqrt{d}),(x_2,y_2/\sqrt{d})\rangle $ and $$\Q(E[N])=\Q(x(E[N]),\sqrt{df(x_1)},\sqrt{df(x_2)}).$$
    We will make use of this equation in several of our results below, namely Prop. \ref{prop-8divfieldm8} and \ref{prop-8divfieldm16}.
\end{remark}

\begin{proposition}\label{prop-8divfieldm16}
Let $d\in \Q^*$ be square-free, and let $E^d\,:\,y^2=x^3 - 11 d^2x + 14d^3$. Then
\begin{itemize}
\item  $\Q(E^d[2]) = \Q(\sqrt{2})$,
%\item  $\Q(E^d[4]) = \Q(\zeta_4,\sqrt{d},\sqrt[4]{2})$.
\item  $\Q(E^d[4])=\Q(\alpha,\sqrt{d})$ such that $f_4(\alpha)=0$, where $f_4(x)=x^8-4x^6+8x^4-4x^2+1$,
 and
	\item $\Q(E^d[8]) = \Q(\beta,\sqrt{d})$ such that $f_8(\beta)=0$, where $f_8(x)=x^{32} + 16x^{31} + 120x^{30} + 560x^{29} + 1848x^{28} + 4784x^{27} + 11000x^{26} + 25344x^{25} + 59844x^{24} + 133856x^{23} +  260768x^{22} + 419392x^{21} + 534920x^{20} + 513536x^{19} + 332032x^{18} + 93856x^{17} - 43548x^{16} - 22112x^{15} + 61056x^{14} + 77728x^{13} + 20768x^{12} - 18304x^{11} + 320x^{10} + 21440x^{9} + 8240x^{8} - 8256x^{7} - 1888x^{6} + 3584x^{5} + 800x^{4} - 1216x^{3} + 320x^{2} + 8
$.
\end{itemize}
In particular, we have $[\Q(E^d[2]):\Q]=2$ and $[\Q(E^d[4]):\Q]=8$, and  we have $[\Q(E^d[8]):\Q]=32$ if and only if $d \in\{\pm 1,\pm2\}$; otherwise $[\Q(E^d[4]):\Q]=16$ and $[\Q(E^d[8]) : \Q] =64$.
\end{proposition}

\begin{proposition}\label{prop-8divfieldm8}
Let $d\in \Q^*$ square-free, and let $E^d\,:\,y^2=x^3 - 4320d^2 x + 96768d^3$. Then
\begin{itemize}
\item  $\Q(E^d[2]) = \Q(\sqrt{2})$.
%\item  $\Q(E^d[4]) = \Q(\zeta_4,\sqrt{d}\sqrt[4]{2},\sqrt{d}\sqrt{2+\sqrt{2}})$.
\item  $\Q(E^d[4])=\Q(\alpha,\sqrt{d}\sqrt[4]{2})$ such that $f_4(\alpha)=0$, where $f_4(x)=x^8+6x^4+1$,
 and
	\item $\Q(E^d[8]) = \Q(\beta,\sqrt{d})$ such that $f_8(\beta)=0$, where $f_8(x)=x^{32} + 16x^{31} + 128x^{30} + 672x^{29} +  2544x^{28} + 7200x^{27} + 15352x^{26} + 24272x^{25} + 26904x^{24} + 17312x^{23} - 304x^{22} - 11984x^{21} - 9672x^{20} - 2720x^{19} - 3592x^{18} - 7552x^{17} -     2798x^{16} + 6224x^{15} + 6368x^{14} - 672x^{13} - 2224x^{12} + 3360x^{11} +     4952x^{10} - 1072x^9 - 4600x^8 - 1120x^7 + 1776x^6 + 752x^5 - 264x^4 -     96x^3 + 24x^2 + 1
$.
\end{itemize}
In particular, we have $[\Q(E^d[2]):\Q]=2$ and $[\Q(E^d[4]):\Q]=16$, and we have $[\Q(E^d[8]):\Q]=32$ if and only if $d \in\{\pm 1,\pm2\}$; otherwise $[\Q(E^d[8]) : \Q] =64$.

\end{proposition}

\begin{proof}
    Let $f(x) = x^3 - 4320x + 96768$. This polynomial factors as 
    \[ f(x) = (x-48)(x + 24 - 36 \sqrt{2})(x + 24 + 36 \sqrt{2}) \]
    over $\Q(\sqrt{2})$. Thus, $\Q(E[2]) = \Q(E^d[2]) = \Q(\sqrt{2})$, which has degree $2$. Next, let $\psi_N(x)$ denote the $N$th division polynomial for $E$. Then 
    \begin{align*}
        \psi_4(x)/2 \psi_2(x) 
        &= (x^2 - 96x - 288)(x^4 + 96x^3 - 12096x^2 + 801792x - 19823616) \\
        &=: g_4(x) \, h_4(x).
    \end{align*}
The quadratic polynomial $g_4(x)$ factors as $g_4(x) = (x - \alpha_1)(x - \alpha_2)$ where $\alpha_1 = 48 + 36 \sqrt{2}$ and $\alpha_2 = 48 - 36 \sqrt{2}$. Let $K_4$ be the splitting field for $h_4(x)$, and suppose $h_4(x)$ factors as $\prod_{j=1}^{4} (x - \beta_j)$, with $\beta_j \in K_4$. By computation, we find that there exist non-square units $u_{\alpha_i} \in \OO_{\Q(\sqrt{2})}$ and $u_{\beta_j} \in \OO_{K_4}$ such that 
\begin{align*}
    f(\alpha_i) = u_{\alpha_i} 3^{6} 2^{8} \sqrt{2} \\
   f(\beta_j) = u_{\beta_j} 3^6 2^8 \gamma^4,
\end{align*}
where $2 \OO_{K_4} = \langle \gamma \rangle^8$.
Again, by computation, we find that
\begin{itemize}
    \item $\sqrt{2} \in K_4$,
    \item $u_{\alpha_i}$ is a square in $\OO_{K_4}$ for $i \in \{1, 2\}$,
    \item $u_{\beta_j}/\sqrt{2}$ is a square in $\OO_{K_4}$ for $1 \leq j \leq 4$.
\end{itemize}
Thus, Remark \ref{rem-howtocomputedivisionfield} shows that  $\Q(x(E^d[4])) = K_4$ and
\[ \Q(E^d[4]) = K_4 \left(\sqrt{d \sqrt{2}}, \sqrt{d u_{\beta_j}} \right) = K_4\left(\sqrt{d} \sqrt[4]{2}\right). \]
Next, we compute the $8$-division field. We find that $\psi_8(x)/2 \psi_4(x) = g_8(x) h_8(x)$, where $g_8(x)$ and $h_8(x)$ are irreducible polynomials of degree $8$ and $16$, respectively. Let $K_{8,g}$ and $K_{8,h}$ be the splitting fields of $g_8(x)$ and $h_8(x)$, which are of degrees $16$ and $32$, respectively.

We find that $K_{8,g}$ is a subfield of $K_{8,h}$, and that for each $\gamma \in K_{8,h}$ where $g_{8}(\gamma) = 0$ or $h_8(\gamma) = 0$, $f(\gamma)$ is a square in $K_{8,h}$. We also note that $\sqrt[4]{2} \in K_{8,h}$, and that the only quadratic fields contained in $K_{8, h}$ are $\Q(\sqrt{2}), \Q(\sqrt{-2})$, and $\Q(\sqrt{-1})$, so
\[ \Q(E^{d}[8]) = 
\begin{cases}
K_{8,h}(\sqrt{d}) & \text{ if } d \neq \pm 1, \pm 2 \\
K_{8,h} & \text{ if } d = \pm 1, \pm 2
\end{cases}, \qquad \text{and} \qquad [\Q(E^{d}[8]) : \Q] = 
\begin{cases}
64 & \text{ if } d \neq \pm 1, \pm 2, \\
32 & \text{ if } d = \pm 1, \pm 2,
\end{cases} \]
as we wanted to prove.
\end{proof}

\end{document}
